\section{Experimental Setup}
\label{sec:setup}

\subsection{Data}
We use the application table of the Home Credit Default Risk dataset~\cite{homecredit2018}, merged with aggregates of
the credit bureau and previous-application tables (Table~\ref{tab:dataset}). The target marks applicants who had
payment difficulties; 8.07\% of the 307{,}511 applicants are positive. The data are split once, stratified by the
target, into 80\% for training and 20\% (61{,}503 applicants) for testing; 20\% of the training part is held out as a
validation set for early stopping and for choosing the cutoffs. Every result is reported on the test set, which is used
for nothing else.

\input{tables/dataset}

Preprocessing follows the base paper~\cite{shreya2025base}. The placeholder used for applicants without employment
(365{,}243 days, almost all of them pensioners) becomes a flag and a missing value. Numeric gaps are filled with the
median and categorical gaps with the most common value, and non-binary numeric columns are clipped to the mean plus or
minus three standard deviations. All these statistics are learned from the training rows only and then applied to the
test rows, so the test data never shape the training data.

\subsection{Models}
The final model uses learning rate 0.03, 35 leaves, maximum depth 6, at least 100 samples per leaf, row and column
subsampling of 0.8, and L1 and L2 regularization of 0.5 and 1.0, with up to 2{,}000 boosting rounds and early stopping
after 150 rounds without improvement on the validation set. No resampling is applied; the three-band router handles the
imbalance through its cutoffs. As baselines, and to compare with the base paper, we train logistic regression (on standardized features), random
forest (200 trees), XGBoost and LightGBM with library defaults otherwise, on SMOTE-balanced training data~\cite{chawla2002smote}, and report
ROC-AUC and the precision, recall and F1 score at each model's F1-optimal threshold. Models are implemented with
LightGBM~\cite{ke2017lightgbm} and scikit-learn~\cite{pedregosa2011sklearn}.

\subsection{Explanation Agreement}
We compare two LIME set-ups with SHAP on 150 test applicants, 50 from each band: LIME on all features with
perturbations drawn from the whole training data, as it is commonly configured, and the local
LIME of Section~\ref{sec:method}. Each is run with five seeds. We report the overlap between the five and ten
most important features of LIME and SHAP; the Spearman rank correlation of the LIME weights with the SHAP values over
SHAP's 25 most important features; the share of SHAP's top five features whose direction LIME reproduces; the
\emph{locality gap}, the rank correlation with the applicant's own SHAP values minus that with another applicant's, which
is near zero for an explanation that is the same for everyone; and the overlap of the top five features across seeds.

\subsection{Recourse Evaluation}
\emph{Applicants.} We draw 300 rejected and 300 Review applicants at random from the test set. Recourse for the
rejected applicants targets Review, and for the Review applicants it targets Approve.

\emph{Methods.} We compare the following methods on the same applicants. Unless stated otherwise, they use the three
loan levers of Section~\ref{sec:recourse} within the same bounds.
\begin{itemize}
\item \emph{Proposed, loan levers only}: the engine with the safety margin, the plausibility check and refinement;
\item \emph{Proposed, with waiting}: the same engine with the waiting lever added;
\item \emph{Ensemble-robust}: the proposed engine with the agreement rule of \eqref{eq:agree}, requiring all five
search models (and, as a variant, four of five), with and without waiting;
\item \emph{Hybrid}: the method used by the application (Section~\ref{sec:robust}), whose first option is evaluated;
\item \emph{Random search}: 300 random points in the same box, with the same margin and plausibility rules, keeping
the cheapest qualifying point;
\item \emph{Baseline engine}: the grid search without the safety margin and the plausibility check, as in the first
version of the system;
\item \emph{DiCE}~\cite{mothilal2020dice}: the model-agnostic random method on the same levers, with its post-hoc
sparsity step, scored through the same feature pipeline; because of its run time it is evaluated on 20 applicants per
tier;
\item \emph{Ablations}: the proposed engine without the margin, without the plausibility check, and with a margin of
0.2.
\end{itemize}

\emph{Measures.} An option is \emph{usable} if one is found, it is plausible and it clears the safety margin; we report
the share of applicants who receive a usable option. We also report the mean cost \eqref{eq:cost} of the options found,
the mean number of levers changed and the run time per applicant.

\emph{Robustness to retraining.} We train ten bootstrap models as described in Section~\ref{sec:robust}. Their test
ROC-AUC is $0.7609 \pm 0.0009$, against 0.7668 for the final model. Models 1--5 are used by the ensemble-robust search;
models 6--10 are never seen by any search and are used only for evaluation. An option \emph{holds} on an evaluation
model if, applied to the applicant, it reaches the promised band with that model's own cutoffs. We report the share of
options that hold on all five evaluation models, and the share of applicants who receive an option that is both usable
and holds on all five.

\subsection{Fairness}
We audit decisions by age band (under 25, 25--35, 35--45, 45--55 and over 55) and by whether the applicant has
children. The demographic parity difference (DPD) is the largest difference in the automatic approval rate between groups,
and the equal opportunity difference~\cite{hardt2016equality} is the same difference among applicants who repaid. To
test whether age needs to be a feature, we retrain the model without the three direct age features, and without them
and the three features derived from age, with the same protocol as the final model, and compare them with the current
model and with a variant that sets the cutoffs separately for each age band.

\subsection{System Measurements}
\emph{Latency.} We submit 100 test applicants through the system's application programming interface (API), after one warm-up applicant, and time
each step and each stage of the decision, on a laptop (Apple silicon, 10 cores, macOS, Python 3.9).

\emph{Understandability.} For 150 test applicants, 50 per band, we rebuild the text an applicant reads, with the
application's own wording, and measure its Flesch reading ease and Flesch--Kincaid grade~\cite{flesch1948readability},
the amount of information shown, the share of displayed features with a written plain-language name, and the number of
technical terms. No study with applicants was conducted; Section~\ref{sec:discussion} discusses this limitation.

\subsection{Reproducibility}
One script rebuilds every model and result from the raw data, and a separate script builds every table and figure in
this paper from the result files. The code, including 279 automated tests, is available at
\url{https://github.com/MKN-Sai-Varun/CreditXplain}.
